\documentclass[10pt,a4paper]{article}
\usepackage[left=2.5cm,right=2.5cm,top=2.5cm,bottom=2.5cm,]{geometry}
\usepackage{amsmath,amssymb,amsthm}
\newtheorem{theorem}{Theorem}[section]
\newtheorem{lemma}[theorem]{Lemma}
\newtheorem{proposition}[theorem]{Proposition}
\usepackage{setspace}
\setlength{\parindent}{0pt}
\onehalfspacing
\usepackage{listings}
\usepackage{xcolor}
\lstset{
    basicstyle=\ttfamily\small,
    keywordstyle=\color{blue},
    commentstyle=\color{green!50!black},
    numbers=left,
    numberstyle=\tiny,
    frame=single,
    breaklines=true
}
\begin{document}
div
\[\vec{n}=\vec{u}\times\vec{v}\]
\[\vec{u}=iu_x+k(\frac{\partial F}{\partial x})u_x\]
\[\vec{v}=jv_y+k(\frac{\partial F}{\partial y})v_y\]
\[\vec{u}\times\vec{v}=(-i\frac{\partial F}{\partial x}-j\frac{\partial F}{\partial y}+k)u_xu_y\]
\[normalize(n)=\frac{-i\frac{\partial F}{\partial x}-j\frac{\partial F}{\partial y}+k}{\sqrt{1+(\frac{\partial F}{\partial x})^2+ (\frac{\partial F}{\partial y})^2}}\]

\[\int\int\int_V\rho(x,y,z,t)dV\]
\[\frac{d}{dt}\int\int\int_V\rho(x,y,z,t)dV\]
\[=\int\int\int_V\frac{\partial \rho}{\partial t}dV=-\int\int_S\rho v\cdot ndS\]
\[\int\int\int_V\frac{\partial\rho}{\partial t}dV=-\int\int\int_V\nabla\cdot(\rho v)dV\]
\[\frac{\partial\rho}{\partial t}=-\nabla\cdot(\rho v)\]
\[\frac{\partial\rho}{\partial t}+\nabla\cdot\rho v=0\]

some relationship with Gauss's law.

\[\int\int_SE\cdot\vec{n}dS=\frac{1}{\varepsilon_0}\int\int\int_V\rho dV\]
\[\int\int\int_V\nabla\cdot EdV=\int\int\int_V\frac{\rho}{\varepsilon_0}dV\]
\[\nabla\cdot E=\frac{\rho}{\varepsilon_0}\]
more style in 2-d version:

\[\int\int_A\nabla\cdot FdA=\int_SF\cdot\vec{n}dS\]
\[F(x,y)=(P(x,y),Q(x,y))\]
\[\nabla\cdot F=\frac{\partial P}{\partial x}+\frac{\partial Q}{\partial y}\]
\[\vec{n}=\begin{pmatrix}0&-1\\1&0\end{pmatrix}\begin{pmatrix}dx\\dy\end{pmatrix}\]
\[c{n}=(dy,-dx)\\[10pt]F\cdot\vec{n}=Pdx-Qdy\]
\[\int\int_A(\frac{\partial P}{\partial x}+\frac{\partial Q}{\partial y})dA=\int(Pdx-Qdy)\]

It's Green's theorem.Just the 2-d version Divergence Theorem.

curl

the curl can be defined:

\[(\nabla\times F)\cdot\vec{n}=\lim_{\Delta S\to0}\frac{1}{\Delta S}\oint F\cdot\vec{t}d\Delta r\]

In Cartesian coordinate:

\[cur F=\nabla\times F=\begin{vmatrix}i&j&k\\\frac{\partial F_x}{\partial x}&\frac{\partial F_y}{\partial y}&\frac{\partial F_z}{\partial z}\\xz&yz&-y^2\end{vmatrix}\]

In cylindrical coordinate:

\[\int_{c_1}D\cdot\vec{t}ds=F(r,\theta-\frac{\Delta\theta}{2},z)\Delta r\]
\[-\int_{c_3}F\cdot\vec{t}ds=-F(r,\theta+\frac{\theta}{2},z)\Delta r\]
\end{document}








